\chapter{Cài đặt và xử lý dữ liệu}
\label{chap:implementation}

\section{Tổ chức chương trình và giao diện}
Mã C++ nằm trong \path{source_code/}. Thư mục khai báo và cài đặt được chia cùng nhóm: mô hình, tính phí, dịch vụ, lưu trữ, ngoại lệ và giao diện. Tệp \code{main.cpp} tạo lớp lưu trữ, dịch vụ và giao diện, sau đó chạy vòng lặp menu. Các lớp nghiệp vụ không cần biết đang được gọi từ menu hay từ chương trình kiểm thử.

\begin{table}[H]
\centering\small
\caption{Các điểm vào của giao diện console}
\label{tab:menu}
\begin{tabularx}{\textwidth}{c l Y}
\toprule
Lựa chọn & Công việc & Tác động dữ liệu\\
\midrule
1 & Nhận xe & Tạo lượt đang gửi nếu lưu thành công.\\
2 & Trả xe & Có thể chỉ xem; chỉ đóng lượt sau xác nhận.\\
3 & Xe đang gửi & Chỉ đọc số xe, sức chứa và các lượt chưa ra.\\
4 & Tìm kiếm & Chỉ đọc theo thẻ hoặc biển số.\\
5 & Lịch sử & Chỉ đọc toàn bộ lượt.\\
6 & Thống kê & Chỉ đọc và tính theo ngày yêu cầu.\\
7 & Đổi sức chứa & Kiểm tra và lưu cấu hình mới.\\
0 & Thoát & Không tạo thêm giao dịch.\\
\bottomrule
\end{tabularx}
\end{table}

Giao diện đọc theo dòng, sau đó kiểm tra số nguyên và miền giá trị. Cách này tránh việc một ký tự sai làm luồng nhập số giữ trạng thái lỗi qua nhiều lần lặp. Khi dữ liệu không hợp lệ, thao tác hiện tại dừng và menu hiện lại. Khi luồng nhập kết thúc, chương trình thoát thay vì tiếp tục hỏi trên một luồng không còn dữ liệu.

Các câu xác nhận yêu cầu nhập \code{y} hoặc \code{Y}. Câu trả lời khác được hiểu là chưa đồng ý. Quy tắc này khiến việc nhấn Enter không tự đóng một lượt gửi. Trên Windows, chương trình thiết lập bảng mã console UTF-8 để thông báo tiếng Việt; phần thiết lập hệ điều hành được giữ trong tệp khởi động, không đưa vào các lớp bãi xe hoặc tính phí.

\section{Luồng xử lý nhận xe}
Sau khi nhận thẻ, biển số và loại xe, giao diện tạo \code{Vehicle}. Nếu biển số không hợp lệ, lỗi được phát hiện trước khi gọi dịch vụ nhận xe. Thời điểm vào được lấy theo chế độ chạy. Dịch vụ sao chép \code{ParkingLot} rồi gọi \code{recordEntry()} trên bản sao để kiểm tra toàn bãi và thêm lượt.

\begin{figure}[H]
\centering
\begin{tikzpicture}[node distance=0.35cm]
\node[reportbox,text width=11.4cm] (input) {Nhập thẻ, biển số, loại xe và lấy thời gian vào};
\node[reportbox,text width=11.4cm,below=of input] (vehicle) {Tạo Vehicle, chuẩn hóa và kiểm tra dữ liệu};
\node[reportbox,text width=11.4cm,below=of vehicle] (copy) {Tạo bản dự kiến của ParkingLot};
\node[reportbox,text width=11.4cm,below=of copy] (check) {Kiểm tra thẻ, biển số đang gửi và sức chứa};
\node[reportbox,text width=11.4cm,below=of check] (save) {Tạo lượt, lưu bản dự kiến qua ParkingRepository};
\node[reportbox,text width=11.4cm,below=of save] (done) {Lưu được: áp dụng vào bộ nhớ, báo nhận xe thành công};
\foreach \a/\b in {input/vehicle,vehicle/copy,copy/check,check/save,save/done}{\draw[reportarrow] (\a)--(\b);}
\end{tikzpicture}

\caption{Luồng nhận xe; chỉ bước cuối thay trạng thái đang chạy}
\label{fig:checkin}
\end{figure}

Trong bản dự kiến, chương trình kiểm tra thẻ và biển số qua các lượt đang hoạt động, sau đó kiểm tra số xe so với sức chứa. Mã lượt chỉ tăng sau khi thêm được đối tượng mới vào danh sách. Dịch vụ chuyển bản dự kiến đến điểm lưu chung; thông báo nhận xe thành công xuất hiện sau khi việc lưu hoàn tất.

Thứ tự này có một chi tiết dễ bỏ sót: bản sao có thể đã chứa lượt mới và mã kế tiếp mới, nhưng bãi đang chạy vẫn chưa thay đổi. Vì vậy, khi không tạo được tệp tạm, cả mã lượt lẫn số xe của bãi thật vẫn giữ nguyên. Khi nhân viên làm lại thao tác sau khi xử lý lỗi, chương trình không để lại một lượt rỗng hoặc tiêu hao mã vì lần nhận trước chưa hoàn tất.

\section{Từ tra thẻ đến xác nhận trả xe}
Phương thức chuẩn bị trả xe lấy lượt đang gửi, kiểm tra thứ tự thời gian và xác định có qua ngày hay không. Nếu cùng ngày, bộ tính phí trả về số tiền. Nếu khác ngày, báo giá đánh dấu cần nhập thủ công và để trống giá trị phí tự động.

\CodeFile{code:prepare-checkout}{Chuẩn bị thông tin trả xe nhưng chưa sửa dữ liệu}{snippets/prepare_checkout.cpp}

Phương thức được khai báo \code{const} vì không đóng lượt hoặc sửa bãi. Kết quả chỉ là dữ liệu để giao diện hỏi nhân viên. Với một lượt cùng ngày, trình tự gọi được mô tả ở Hình~\ref{fig:checkout-sequence}. Hai khoảng làm việc của nhân viên nằm giữa bước xem báo giá và bước xác nhận.

\begin{figure}[H]
\centering
\begin{tikzpicture}[font=\small]
\foreach \x/\name in {0/Giao diện,3/Dịch vụ,6/Bãi xe,9/Tính phí,12/Lưu trữ}{
\node[reportbox,text width=2cm,minimum height=0.7cm] at (\x,0) {\name};
\draw[dashed,draw=rulegray] (\x,-0.5)--(\x,-7.2);}
\draw[reportarrow] (0,-0.9)--node[above,fill=white,inner sep=1pt]{Tra thẻ}(3,-0.9);
\draw[reportarrow] (3,-1.7)--node[above,fill=white,inner sep=1pt]{Tìm lượt}(6,-1.7);
\draw[reportarrow] (3,-2.5)--node[above,fill=white,inner sep=1pt]{Tính phí cùng ngày}(9,-2.5);
\draw[reportarrow,dashed] (3,-3.3)--node[above,fill=white,inner sep=1pt]{Báo giá}(0,-3.3);
\node[align=left,fill=white,inner sep=2pt] at (1.4,-4.1) {Đối chiếu biển số\\Xác nhận thu tiền};
\draw[reportarrow] (0,-4.9)--node[above,fill=white,inner sep=1pt]{Xác nhận}(3,-4.9);
\draw[reportarrow] (3,-5.7)--node[above,fill=white,inner sep=1pt]{Lưu trạng thái dự kiến}(12,-5.7);
\draw[reportarrow,dashed] (12,-6.3)--node[above,fill=white,inner sep=1pt]{Lưu thành công}(3,-6.3);
\draw[reportarrow,dashed] (3,-6.9)--node[above,fill=white,inner sep=1pt]{Cho xe ra}(0,-6.9);
\end{tikzpicture}

\caption{Trình tự rút gọn của một lượt trả xe cùng ngày}
\label{fig:checkout-sequence}
\end{figure}

Khi gọi \code{confirmCheckout()}, dịch vụ tra lại theo thẻ và thời gian ra, rồi so sánh mã lượt, biển số và giờ vào. Phí cùng ngày được lấy từ kết quả tính lại; số tiền do đối tượng báo giá mang đến không được tin ngay. Với lượt qua ngày, số tiền và lý do được kiểm tra trước khi ghi vào lượt dự kiến.

Nhân viên xác nhận thu tiền trước khi dịch vụ lưu dữ liệu, nhưng thông báo cho xe ra chỉ xuất hiện sau khi lưu thành công. Nếu bước lưu thất bại, màn hình phải báo lỗi, xe vẫn ở trạng thái đang gửi trong phần mềm. Bản console chưa có quy trình đối soát tiền mặt cho trường hợp nhân viên đã nhận tiền nhưng lưu bị lỗi; đây là một điểm cần thiết kế thêm trước khi đưa vào vận hành thực tế.

\section{Tệp dữ liệu và cách khôi phục}
Lớp lưu trữ sử dụng tệp văn bản để nhóm dễ quan sát dữ liệu khi kiểm tra. Dòng đầu có phiên bản định dạng, sức chứa và mã lượt kế tiếp. Các dòng sau chứa từng lượt gửi. Cấu trúc này đủ cho ứng dụng một phiên, không yêu cầu cài thêm một hệ quản trị cơ sở dữ liệu.

\begin{table}[H]
\centering\small
\caption{Các trường của một lượt trong tệp dữ liệu}
\label{tab:file-fields}
\begin{tabularx}{\textwidth}{l Y}
\toprule
Trường & Cách hiểu khi đọc lại\\
\midrule
Mã lượt & Số nguyên dương, không trùng và nhỏ hơn mã kế tiếp.\\
Mã thẻ, biển số & Chuỗi được đọc rồi chuẩn hóa, kiểm tra hợp lệ.\\
Loại xe & 1 cho xe máy, 2 cho xe máy điện.\\
Thời gian vào & Chuỗi ngày giờ có độ chính xác đến phút.\\
Thời gian ra & Ngày giờ hợp lệ, hoặc \code{-} nếu đang gửi.\\
Phí đã thu & Số nguyên không âm; lượt đang gửi có giá trị 0.\\
Cờ thủ công & 0 hoặc 1; phải phù hợp với việc qua ngày.\\
Lý do & Chuỗi một dòng; bắt buộc khi thu thủ công.\\
\bottomrule
\end{tabularx}
\end{table}

Các chuỗi được ghi bằng \code{std::quoted}. Nhờ đó, lý do có dấu cách, dấu nháy hoặc dấu gạch chéo vẫn đọc lại đúng. Một bản dữ liệu minh họa có thể viết như sau:

\begin{Code}
PARKING_DATA_V1 100 3
1 "C001" "29A112345" 1 "2026-10-26 08:00" "2026-10-26 17:00" 3000 0 ""
2 "C002" "29M99999" 2 "2026-10-26 18:00" "-" 0 0 ""
\end{Code}

Không thể chỉ đọc đủ trường rồi coi dữ liệu đã hợp lệ. Lớp lưu trữ kiểm tra phiên bản, dữ liệu dư, kiểu xe, ngày giờ, phí và cờ thủ công. Sau đó, khi tạo lại bãi, các mã lượt, thẻ và biển số đang gửi được kiểm tra trùng; số xe phải không vượt sức chứa. Phí của lượt cùng ngày cũng được tính lại để đối chiếu với giá trị lưu.

Nếu chưa có tệp chính và chưa có bản sao lưu, chương trình hỏi sức chứa để tạo bãi mới. Nếu có tệp nhưng nội dung sai, chương trình dừng khởi động với thông báo lỗi. Hai trường hợp phải được tách rõ: tệp hỏng không có nghĩa là chưa từng có dữ liệu và không cho phép ghi một bãi trống đè lên lịch sử cũ.

\begin{figure}[H]
\centering
\begin{tikzpicture}[node distance=0.35cm]
\node[reportbox,text width=11.6cm] (write) {Ghi trạng thái dự kiến vào tệp .tmp};
\node[reportbox,text width=11.6cm,below=of write] (flush) {Kiểm tra ghi, đẩy dữ liệu và đóng tệp};
\node[reportbox,text width=11.6cm,below=of flush] (backup) {Nếu có tệp chính: đổi tên thành .bak};
\node[reportbox,text width=11.6cm,below=of backup] (replace) {Đổi .tmp thành tệp chính};
\node[reportbox,text width=11.6cm,below=of replace] (finish) {Thành công: áp dụng trạng thái; dọn bản sao lưu};
\foreach \a/\b in {write/flush,flush/backup,backup/replace,replace/finish}{\draw[reportarrow] (\a)--(\b);}
\node[below=0.25cm of finish,text width=12cm,align=center,font=\small] {Nếu thay thế thất bại: cố khôi phục .bak.\\Nếu thiếu tệp chính khi khởi động: có thể đọc .bak.};
\end{tikzpicture}

\caption{Các bước thay thế tệp dữ liệu và vị trí bản sao lưu}
\label{fig:storage}
\end{figure}

Khi lưu, chương trình ghi toàn bộ trạng thái dự kiến vào \code{.tmp} và kiểm tra luồng sau khi ghi, đẩy dữ liệu, đóng tệp. Nếu có tệp chính, nó được chuyển sang \code{.bak} trước khi tệp tạm được đưa vào vị trí chính. Khi thay thế thất bại, chương trình cố đưa bản cũ trở lại. Sau khi tệp chính mới đã hoàn tất, lỗi dọn bản sao lưu không làm đảo kết quả giao dịch.

Nếu tiến trình dừng đúng lúc tệp chính đã đổi tên nhưng tệp mới chưa vào vị trí, lần mở sau có thể đọc \code{.bak} khi thiếu tệp chính. Cơ chế này hỗ trợ tình huống gián đoạn đã mô phỏng, nhưng không phải cam kết phục hồi mọi lỗi mất điện. Chương trình chưa có khóa giữa nhiều tiến trình, nhật ký giao dịch hoặc cơ chế ép dữ liệu xuống thiết bị lưu trữ như một cơ sở dữ liệu.

\section{Ngày giờ, tiền và thống kê}
\code{DateTime} lưu năm, tháng, ngày, giờ và phút dưới dạng các giá trị số. Hàm tạo kiểm tra miền giá trị và số ngày của tháng, bao gồm quy tắc năm nhuận. Hàm phân tích chuỗi yêu cầu đúng dạng \code{YYYY-MM-DD HH:MM}; không chấp nhận một ngày sai rồi tự chuyển sang tháng sau.

\CodeFile{code:datetime}{Kiểm tra ngày giờ ngay khi tạo đối tượng}{snippets/datetime.cpp}

So sánh thời điểm dựa trên thứ tự năm, tháng, ngày, giờ, phút. Vì dữ liệu trong bài toán là giờ địa phương trên một máy, cách biểu diễn này thuận tiện cho việc so cùng ngày và so với 19:00. Chương trình chưa xử lý việc thay đổi múi giờ hoặc điều chỉnh đồng hồ giữa các cổng. Nếu đồng hồ dẫn đến giờ ra trước giờ vào, thao tác bị từ chối để nhân viên kiểm tra lại.

Số tiền dùng \code{long long}; các mức phí tự động được trả về bằng số nguyên. Khi hiển thị, giao diện thêm dấu chấm phân cách hàng nghìn, nhưng tệp vẫn lưu số nguyên không có định dạng trang trí. Lý do thu thủ công được giới hạn trên một dòng, tối đa 500 byte UTF-8, để một bản ghi không bị tách thành nhiều dòng.

\section{Build và chạy chương trình}
Chương trình dùng C++17 và thư viện chuẩn, đã build trên Windows bằng GCC 16.1.0. Chế độ thường dùng giờ máy; demo cho nhập giờ và dùng dữ liệu riêng. Từ \path{source_code/}, chạy:

\begin{Code}
powershell -ExecutionPolicy Bypass -File scripts/build.ps1
.\build\parking.exe --demo
powershell -ExecutionPolicy Bypass -File scripts/test.ps1
\end{Code}

